\documentclass[11pt]{article}
\usepackage[T1]{fontenc}
\usepackage[utf8]{inputenc}
\usepackage{lmodern,amsmath,amssymb,amsthm,mathtools}
\usepackage[margin=1in]{geometry}
\usepackage{microtype,booktabs,longtable,array,enumitem}
\usepackage[hidelinks]{hyperref}
\setlength{\emergencystretch}{2em}
\setlist{itemsep=2pt,topsep=5pt}
\newtheorem{theorem}{Theorem}[section]
\newtheorem{lemma}[theorem]{Lemma}
\newtheorem{proposition}[theorem]{Proposition}
\newtheorem{corollary}[theorem]{Corollary}
\theoremstyle{definition}
\newtheorem{definition}[theorem]{Definition}
\newtheorem{example}[theorem]{Example}
\newtheorem{remark}[theorem]{Remark}
\newcommand{\R}{\mathbb R}
\newcommand{\Q}{\mathbb Q}
\newcommand{\Z}{\mathbb Z}
\newcommand{\Ssym}{\mathbb S}
\newcommand{\CP}{\mathcal{CP}}
\newcommand{\COP}{\mathcal{COP}}
\newcommand{\GG}{\mathcal G_{\Q}}
\newcommand{\AC}{\mathfrak A}
\newcommand{\cone}{\operatorname{cone}}
\newcommand{\lin}{\operatorname{span}}
\newcommand{\ri}{\operatorname{ri}}
\newcommand{\tr}{\operatorname{tr}}
\newcommand{\diag}{\operatorname{diag}}
\newcommand{\cl}{\operatorname{cl}}
\newcommand{\rank}{\operatorname{rank}}
\newcommand{\dist}{\operatorname{dist}}
\newcommand{\Min}{\operatorname{Min}}
\newcommand{\htb}{\operatorname{ht}_2}
\newcommand{\ip}[2]{\langle #1,#2\rangle}
\newcommand{\norm}[1]{\lVert #1\rVert}
\newcommand{\trans}{^{\mathsf T}}
\numberwithin{equation}{section}
\title{Arithmetic carriers of constrained Gram representations\\[0.4em]
\large Rational completion, sharp moment corrections, and the limits of local traversal}
\author{Manuscript II --- mathematical research draft}
\date{15 November 2026}
\begin{document}
\maketitle
\begin{abstract}
Let the columns of a Gram factor belong to a finite union of rational polyhedral cones, and impose rational affine observations on the Gram matrix. We prove that the closure of the viable rational column directions is itself a finite union of closed rational polyhedral cones. This \emph{arithmetic carrier} is intrinsic and is computed by at most $n(n+1)/2$ strict matrix-span reductions. Algebraic quadratic exposers, stationarity on minimal vector faces, and rational coefficient kernels give an exact factor-or-obstruction algorithm, with polynomial bit complexity in every fixed matrix order. The carrier characterizes measurement-preserving rational regularization, including each prescribed rational direction, and yields descent of polyhedrally constrained Gram representations from totally real number fields. For bounded moment matrices we determine the exact correction for every tensor grid: it is the square sine of half the largest angular gap. Rational stereographic grids attain quadratic accuracy, whereas the geometrically optimal grid is not rational once it has at least four intervals. Finally, in every dimension at least four a fixed nonsimplicial cone has interior rational targets with constant-width factors but unique shadow walks of exponential binary length. Their carrier and observation condition remain fixed. Thus genuine arithmetic faces control certification, whereas an interior perfect-form subdivision can impose an independent local traversal cost.
\end{abstract}
\noindent\textbf{Scope of this draft.} The arguments below are self-contained except for explicitly cited real-algebraic algorithmic primitives and classical convex separation. The accompanying audit does not certify historical priority or a journal percentile. The exact programs verify specified certificates, not the general quantifier-elimination algorithm or a formal proof of the universal theorems.
\par\smallskip
\noindent\textbf{Degenerate-case conventions.} Every vector domain is the union of at least one cone, so it contains zero. The lattice conditioning statement assumes positive span dimension. Observation coercivity and the observed-diagonal corollary concern a nonzero Gram cone; for the zero cone the only possible completion is the zero matrix, decided directly. These conventions apply throughout.
\tableofcontents

\section{The distinction that governs the problem}
A rational matrix, a rational Gram factor, and a rational point in a prescribed affine fiber are different objects. For a closed vector domain $P\subseteq\R^n$ that is a union of cones, write
\begin{align}
 C(P)&=\cone\{xx\trans:x\in P\},\\
 \GG(P)&=\left\{\sum_{j=1}^s u_ju_j\trans:s<\infty,\ u_j\in P\cap\Q^n\right\}.
\end{align}
The empty sum is permitted. Throughout the main theorem, $P$ is a finite union of closed rational polyhedral cones; components may have lineality or smaller dimension. In particular $P$ need not be convex. Let
\begin{equation}\label{eq:fiber}
 H=\{X\in\Ssym^n:\mathcal L(X)=b\}
\end{equation}
be rational affine, and put $d=n(n+1)/2$. For a convex cone $K$, $C(K)$ is the generalized completely positive cone and its dual consists of the quadratic forms nonnegative on $K$.

Our main object is a vector-domain invariant, not a pivot potential. Define the \emph{viable rational directions}
\begin{equation}\label{eq:viable}
 \mathcal D_{\Q}(P,H)=\{u\in P\cap\Q^n:\exists X\in\GG(P)\cap H,\ \exists\tau\in\Q_{>0},\ X-\tau uu\trans\in\GG(P)\}.
\end{equation}
The coefficient $\tau$ cannot be omitted: an arbitrarily long prescribed column cannot fit a fixed trace. Define
\begin{equation}\label{eq:carrierdef}
 \AC(P,H)=\cl\mathcal D_{\Q}(P,H).
\end{equation}
If the rational fiber is empty, its carrier is empty. If the only feasible matrix is zero, the carrier is $\{0\}$, not the empty set.

\begin{theorem}[Canonical arithmetic carrier]\label{thm:main}
For the preceding rational data there is a finite algorithm which either proves $\GG(P)\cap H=\varnothing$, or computes $\AC=\AC(P,H)$ as a finite union of closed rational polyhedral subcones of $P$. It uses at most $d$ proper matrix-span reductions. In the nonempty case,
\begin{align}
 \mathcal D_{\Q}(P,H)&=\AC\cap\Q^n,\label{eq:directions}\\
 \GG(P)\cap H&=\GG(\AC)\cap H,\label{eq:exactfiber}\\
 \cl(\GG(P)\cap H)&=C(\AC)\cap H,\label{eq:closure}\\
 H\cap\ri C(\AC)&\ne\varnothing.\label{eq:slater}
\end{align}
The set $\AC$, not merely the matrix slice in \eqref{eq:closure}, is independent of all choices of exposing forms. There is a rational $X_0\in H\cap\ri C(\AC)$ which witnesses every direction in \eqref{eq:directions}, with $\tau$ allowed to depend on $u$. Such a witness has a full-span reserve representation with at most $2d$ weighted rational atoms. Every feasible rational matrix admits at most $d$ weighted rational atoms, hence at most $4d$ rational Gram columns, whenever it belongs to $\GG(P)$.
\end{theorem}

The last width assertion is an existence and output-size statement; it is not a claim of a deterministic polynomial-time four-square routine. Section~\ref{sec:complexity} distinguishes that width from the polynomial-time output construction. The theorem does not assert that every rational point of the real fiber factors rationally. Rather, it computes the exact arithmetic obstruction.

Classical facial reduction supplies successive exposing functionals~\cite{BW,LMT}. The additional step here is at the level of \emph{column directions}: an algebraic zero restriction is replaced by its exact rational coefficient kernel before proceeding. Rational interior factorization~\cite{DSV17}, rational-point algorithms for convex semialgebraic sets~\cite{SEDZ}, and generalized perfect-matrix theory~\cite{OS26} are antecedents, not claims of this paper. In particular, deciding whether a rational matrix point exists in $C(P)\cap H$ does not decide whether that point has permitted rational columns.

The proof also explains a limitation of unification by a single facial depth. A perfect-form normal fan subdivides the \emph{interior} of a matrix cone. Its cells are not genuine faces of that cone. Section~\ref{sec:paths} constructs exponentially many consecutive cells while the genuine arithmetic carrier is unchanged. The corrected principle is therefore a certification theorem together with a precise obstruction to using it as a polynomial local-walk bound.

\section{Rank-one reserves and arithmetic zero domains}\label{sec:zeros}
We use the Frobenius pairing on symmetric matrices. A rational weighted representation means $X=\sum_j\lambda_j u_ju_j\trans$ with $\lambda_j\in\Q_{\ge0}$ and $u_j\in P\cap\Q^n$.

\begin{lemma}[Weights, closedness, and a rational reserve]\label{lem:reserve}
Rational weighted representations are equivalent to membership in $\GG(P)$. The real cone $C(P)$ is closed, its linear span $V$ is rational, and there are rational permitted atoms $G_i=a_ia_i\trans$, $1\le i\le k=\dim V$, forming a basis of $V$. Their sum $B=\sum_iG_i$ belongs to $\ri C(P)$.
\end{lemma}
\begin{proof}
For $\lambda=a/b>0$ in lowest terms, Lagrange's four-square theorem gives $ab=\sum_{\ell=1}^4s_\ell^2$ with nonnegative integers $s_\ell$. Thus
\[
 \lambda uu\trans=\sum_{\ell=1}^4((s_\ell/b)u)((s_\ell/b)u)\trans.
\]
Positive rescaling preserves every component cone. The reverse implication is immediate.

If $P\ne\{0\}$, the normalized atoms $xx\trans$ with $x\in P$, $\norm{x}_2=1$, form a compact set. Its convex hull is exactly the trace-one base of $C(P)$ and is compact. Since trace equals one on that base, its cone is closed. The zero case is immediate.

For each component choose rational cone vectors $a_1,\ldots,a_r$ spanning its vector span. Their squares and the squares of $a_i+a_j$ span the symmetric matrices supported on that span. All these vectors remain in the component. Taking the union over components and selecting an atom basis gives $G_1,\ldots,G_k$. The simplicial cone generated by that basis is full-dimensional in $V$; its positive sum has a relative neighborhood in that cone. Hence $B\in\ri C(P)$.
\end{proof}

\begin{lemma}[Rational relative-interior factorization]\label{lem:interior}
Every rational $X\in\ri C(P)$ belongs to $\GG(P)$. It has a rational weighted representation using at most $2k$ atoms that includes a strictly positive coefficient on every atom of a basis of $V$. Without the reserve requirement, at most $k$ weighted atoms suffice.
\end{lemma}
\begin{proof}
Choose rational $\tau>0$ sufficiently small that $X-\tau B\in C(P)$. Conic Carath\'eodory gives
\[
 X-\tau B=\sum_{j=1}^{k}x_jx_j\trans,
\]
allowing zero summands and absorbing real nonnegative weights into the vectors. Approximate each $x_j$ by rational $r_j$ in the same rational polyhedral component. In the rational basis $(G_i)$ the residual
\[
 X-\sum_jr_jr_j\trans=\sum_i e_iG_i
\]
has rational coordinates $e_i$ converging to $\tau$. Sufficiently accurate approximations make every $e_i>0$. This gives the asserted reserve with at most $2k$ atoms. A rational linear dependence among more than $k$ atoms can be eliminated by moving the rational coefficient vector to the first nonnegative boundary along a rational null vector. Each step removes an atom and preserves rational nonnegative coefficients. Lemma~\ref{lem:reserve} then converts weights to Gram columns.
\end{proof}

\begin{lemma}[Facewise stationarity]\label{lem:stationarity}
Let $Q$ be a closed polyhedral vector cone, let $W$ be nonnegative on $Q$, and suppose $x\in Q$ satisfies $x\trans Wx=0$. If $F$ is the minimal face of $Q$ containing $x$, then the restriction of $W$ to $U=\lin F$ is positive semidefinite and annihilates $x$.
\end{lemma}
\begin{proof}
Since $x\in\ri F$, for every $v\in U$ the points $x+tv$ lie in $F$ for sufficiently small positive and negative $t$. The polynomial
\[
 (x+tv)\trans W(x+tv)=2t\,x\trans Wv+t^2v\trans Wv
\]
is nonnegative on a two-sided interval about zero. Its linear coefficient vanishes and its quadratic coefficient is nonnegative. This proves both assertions. The proof includes lineality and $x=0$.
\end{proof}

\begin{lemma}[Exact rational zero domains]\label{lem:zeros}
Let $W$ be a real algebraic symmetric matrix nonnegative on $P$. For every face $F$ of every component, choose a rational basis matrix $B_F$ for $\lin F$. In a common number field $\Q(\theta)$, using an irreducible power basis, expand
\begin{equation}\label{eq:coeff}
 B_F\trans WB_F=\sum_{a=0}^{e-1}\theta^a M_{F,a},\qquad M_{F,a}\in\Q^{r_F\times r_F}.
\end{equation}
Set
\begin{equation}\label{eq:DF}
 D_F=F\cap\left\{B_Fz:M_{F,a}z=0\ \text{for every }a\right\}.
\end{equation}
Then
\[
 \cl\{u\in P\cap\Q^n:u\trans Wu=0\}=\bigcup_F D_F.
\]
In particular this closure is a finite union of closed rational polyhedral cones, contained in the real zero set.
\end{lemma}
\begin{proof}
Take a rational zero $u$ and its minimal component face $F$. Its coordinates $z$ in $B_F$ are rational. Lemma~\ref{lem:stationarity} gives $(B_F\trans WB_F)z=0$. Linear independence of the power basis in \eqref{eq:coeff} gives $M_{F,a}z=0$ for every $a$, so $u\in D_F$.

Conversely every point of $D_F$ has zero quadratic value. Each $D_F$ is rational polyhedral, so its rational points are dense in it. The finite union is closed. The two inclusions now give the identity. No positive-semidefiniteness test for arbitrary $F$ is needed in the definition: stationarity is used only to prove that every rational zero is captured.
\end{proof}

\section{Proof and canonicity of the carrier theorem}\label{sec:carrierproof}
The algorithm maintains a finite union of rational polyhedral cones, initially $P$.

\begin{quote}
\textbf{Arithmetic carrier algorithm.}
Compute $V=\lin C(P)$ and a rational reserve $B$ as in Lemma~\ref{lem:reserve}. Decide whether $C(P)\cap H$ is empty. If empty, return EMPTY. If $H\cap\ri C(P)\ne\varnothing$, return $P$. Otherwise select an algebraic $W$ that is nonnegative on $P$, nonzero on $V$, zero on $H\cap V$, and normalized by $\ip{W}{B}=1$. Replace $P$ by the union in Lemma~\ref{lem:zeros}, and repeat.
\end{quote}
Every real test and selection in this description has an exact real-algebraic implementation; its cost is included in Section~\ref{sec:complexity}.

\begin{lemma}[Availability of an exposer]\label{lem:exposer}
If $C(P)\cap H\ne\varnothing$ but $\ri C(P)\cap H=\varnothing$, the normalized algebraic exposer required by the algorithm exists.
\end{lemma}
\begin{proof}
Work in $V$. Separate the relatively open convex set $\ri C(P)$ from the affine space $H\cap V$. The separating functional is constant on that affine space. Scaling points of the cone and using a feasible point shows that this constant is zero and that the functional is nonnegative on the cone. It is nonzero on $V$. Since $B\in\ri C(P)$, its value at $B$ is strictly positive, so normalization is possible.

The requirements are a first-order real formula over $\Q$: linear equations express vanishing on the rational affine space and the normalization, while a universal vector block expresses nonnegativity on the finite union of rational polyhedral cones. A nonempty semialgebraic set defined over $\Q$ has a point with real algebraic coordinates. Thus an algebraic $W$ can be selected. The zero-dimensional case has already terminated, since feasibility there implies relative strict feasibility.
\end{proof}

\begin{proof}[Proof of Theorem~\ref{thm:main}]
At one reduction let $P'$ be the new domain. It lies in $P$, and $C(P')\subseteq C(P)\cap W^\perp$. Because $W$ is nonzero on $V$, the matrix span strictly decreases.

Every rational feasible Gram factorization survives \emph{column by column}. Indeed, if $X=\sum_j u_ju_j\trans\in H$, then
\[
 0=\ip{W}{X}=\sum_j u_j\trans Wu_j
\]
is a sum of nonnegative quantities. Every $u_j$ is a rational zero and belongs to $P'$ by Lemma~\ref{lem:zeros}. Conversely $P'\subseteq P$. Thus the rational fiber is preserved exactly at each step. There can be at most $d$ proper reductions. A later empty real-feasibility test proves that the original rational fiber was empty.

Suppose the procedure terminates at a nonempty $P_*$. The affine space $H\cap\lin C(P_*)$ is rational, and its intersection with $\ri C(P_*)$ is nonempty and relatively open. It therefore contains a rational point $X_0$. Lemma~\ref{lem:interior} factors $X_0$ rationally and gives a full-span reserve.

For any $X\in C(P_*)\cap H$, the segment $(1-t)X+tX_0$ lies in $H\cap\ri C(P_*)$ for $0<t\le1$. Rational points in the rational affine space can approximate these segment points while remaining in that relatively open set. Lemma~\ref{lem:interior} factors all such rational approximants. This proves \eqref{eq:closure}; the reverse closure inclusion follows from closedness. Exact preservation gives \eqref{eq:exactfiber}.

It remains to prove that $P_*$ itself is intrinsic. If $u\in\mathcal D_{\Q}(P,H)$, decompose both the residual in \eqref{eq:viable} and its positive rational weight into rational squares. The preceding columnwise argument shows that every positive rescaling of $u$ occurring in this factorization survives every step. Hence $u\in P_*$. Conversely let $u\in P_*\cap\Q^n$. For sufficiently small rational $\tau>0$, $X_0-\tau uu\trans$ remains in $\ri C(P_*)$. It is rational, and Lemma~\ref{lem:interior} factors it. Thus $u$ is viable. We have proved
\[
 \mathcal D_{\Q}(P,H)=P_*\cap\Q^n.
\]
Rational points are dense in the finite union $P_*$, which is closed. Therefore $P_*=\cl\mathcal D_{\Q}(P,H)=\AC(P,H)$, independent of the exposers. The width and reserve assertions follow from Lemmas~\ref{lem:reserve} and~\ref{lem:interior}.
\end{proof}

\begin{corollary}[A rank-one arithmetic alternative]\label{cor:insertion}
Let $u\in P\cap\Q^n$. Either $u$ is viable, with a rational insertion certificate \eqref{eq:viable}, or the carrier algorithm produces an empty-fiber certificate or a first exposer strictly positive at $u$. In the latter case the exposer is nonnegative on the \emph{current} domain and annihilates its feasible matrix slice. The preceding reduction chain is part of the certificate.
\end{corollary}
\begin{proof}
A rational zero always survives an update. Thus a rational direction first removed from a nonempty chain has strictly positive value at that step. The converse follows by pairing a hypothetical surviving rational factorization with that exposer.
\end{proof}

A full-span reserve makes insertion quantitative. Write
\[
 X_0=Y+\sum_{i=1}^k e_iG_i,\quad e_i>0,\quad Y\in\GG(\AC),\qquad uu\trans=\sum_{i=1}^ka_iG_i.
\]
Then any rational
\begin{equation}\label{eq:tauinsert}
 0<\tau\le\frac{\min_i e_i}{2(1+\max_i|a_i|)}
\end{equation}
works, since every $e_i-\tau a_i$ remains positive. Cramer's rule bounds the extra binary height polynomially in that of the stored reserve and $u$, for fixed $n$. The same construction inserts any finite list of viable rational directions simultaneously, after reducing their coefficients sufficiently.


\section{Exact observations and every rational regularization direction}\label{sec:directions}
Let $A\in C(P)\cap H$ be rational and let $E$ be the direction space of $H$.

\begin{theorem}[Structured rational regularization]\label{thm:regularize}
There exists $D\in E\cap\Ssym^n(\Q)$ such that $A+tD\in\GG(P)$ for every rational $0<t\le1$ if and only if
\[
 A\in C(\AC(P,H))\cap H.
\]
When the carrier is empty the condition is false. When it holds one may take $D=X_0-A$, where $X_0$ is the rational reserve completion of Theorem~\ref{thm:main}.
\end{theorem}
\begin{proof}
Necessity is \eqref{eq:closure}. For sufficiency, every positive point on the segment from $A$ to $X_0$ belongs to $\ri C(\AC)$, is rational for rational $t$, and lies in $H$. Apply Lemma~\ref{lem:interior}.
\end{proof}

This criterion decides the existence of a fixed perturbation direction, rather than merely a sequence of unrelated rational approximants. An individual prescribed direction admits an exact test as well.

\begin{theorem}[Prescribed-direction criterion]\label{thm:individual}
Let $D\ne0$ be rational and apply the carrier algorithm to $H_D=A+\R D$. If the output is nonempty, define the closed interval
\[
 I_D=\{t\in\R:A+tD\in C(\AC(P,H_D))\}.
\]
Then $A+tD$ has permitted rational Gram columns for every sufficiently small positive rational $t$ if and only if $0\in I_D$ and $I_D$ contains a positive point. For $D=0$, the condition is exactly $A\in\GG(P)$.
\end{theorem}
\begin{proof}
The closure of the rational feasible parameter set is $I_D$, by \eqref{eq:closure} and the homeomorphism between the line and $\R$. The parameter set is closed under rational convex combinations, by Lemma~\ref{lem:reserve}. Every rational parameter in the interior of $I_D$ can be bracketed by two rational feasible parameters; interpolation with a rational coefficient proves its feasibility. The stated interval condition puts all sufficiently small positive parameters in that interior. Necessity follows by taking limits and one positive feasible parameter.
\end{proof}

\begin{corollary}[Completion and coefficient-constrained squares]\label{cor:applications}
The carrier algorithm decides rational Gram completion with arbitrary rational prescribed entries and rational polyhedral restrictions on each factor column. It also decides rational sum-of-squares representation with rational affine conditions on the resulting polynomial and with each summand's coefficient vector constrained to a prescribed finite union of rational polyhedral cones. The same carrier gives the exact regularization and insertion criteria in both settings.
\end{corollary}
\begin{proof}
Matrix entries are linear measurements. For a fixed monomial vector $v(x)$, the coefficient vector of $v(x)\trans Xv(x)$ is a rational linear function of $X$. A decomposition $X=\sum u_ju_j\trans$ is exactly the polynomial representation $\sum (u_j\trans v(x))^2$. Apply the main theorem to those measurements and column domains.
\end{proof}

The unstructured sum-of-squares case is classical~\cite{SEDZ,Scheiderer}; the assertion here retains arbitrary polyhedral constraints on the coefficients of the individual squares. Likewise, for nonnegative columns a prescribed zero $X_{ij}=0$ forces $u_i u_j=0$ columnwise. A graph zero pattern can therefore be incorporated by a union of rational coordinate cones on the permitted cliques. This implication is not used for signed columns, where cancellation may occur.

\section{Descent from totally real fields}\label{sec:descent}
The relevant descent is not obtained by averaging admissible columns: a field embedding can violate their cone inequalities.

\begin{theorem}[Polyhedrally constrained totally real descent]\label{thm:descent}
A rational affine Gram problem $(P,H)$ has a rational permitted Gram solution if and only if it has a solution whose columns belong to a finite totally real number field. The columns are unweighted Gram columns; alternatively, one may allow weights that are totally positive in that field. Positivity in only the selected embedding is not sufficient for the weighted version.
\end{theorem}
\begin{proof}
Only the reverse implication requires proof. Let $E$ be totally real and let $X=\sum_jx_jx_j\trans\in H$ be a surviving $E$-valued factorization at a step of the carrier algorithm. For each $j$ choose the minimal current polyhedral face $F_j$ containing $x_j$. The selected exposer $W$ has zero value at every $x_j$, and is positive semidefinite on $\lin F_j$ by Lemma~\ref{lem:stationarity}.

Every embedding $\sigma:E\to\R$ sends $x_j$ into the same rational linear span, although it need not preserve the face inequalities. The matrix
\[
 \overline X=\frac1{[E:\Q]}\sum_\sigma\sum_j\sigma(x_j)\sigma(x_j)\trans
\]
is rational, belongs to $H$, and lies in the current rational matrix span. Hence $\ip{W}{\overline X}=0$. Every term is nonnegative by the positive-semidefinite restrictions just established. Consequently those restrictions annihilate every conjugate $\sigma(x_j)$.

Expand $x_j$ in a rational basis of $E$. The embedding matrix is nonsingular, so the real span of the conjugates equals the real span of its rational coefficient vectors. The restricted $W$ annihilates each coefficient vector. Expanding $W$ in its own irreducible number-field basis then shows that all rational coefficient-kernel equations \eqref{eq:DF} annihilate these vectors. Therefore the original $x_j$ lies in the replacement domain. The whole factorization survives every step; no real-feasibility test becomes empty. At the nonempty terminal stage Lemma~\ref{lem:interior} and rational affine density give a rational solution.

With totally positive weights, insert their positive conjugates in the trace sum. Each term remains nonnegative, and the same conclusion holds. The forward implication uses the field $\Q$.
\end{proof}

Unconstrained totally real descent is due to Hillar, Quarez and the trace-form treatment of Scheiderer~\cite{Hillar,Scheiderer}. The preceding proof uses that trace argument only after facewise stationarity makes conjugation harmless on the appropriate linear spans.

\begin{example}[Real feasibility without a rational completion]\label{ex:arithmetic}
Consider the rational affine line
\[
 A(t)=\diag\left(\begin{pmatrix}t&2\\2&2t\end{pmatrix},\begin{pmatrix}2&t\\t&1\end{pmatrix}\right).
\]
Positive semidefiniteness forces $t\ge0$, $2t^2-4\ge0$, and $2-t^2\ge0$, hence $t=\sqrt2$. At that parameter the nonnegative columns
\[
 (2^{1/4},2^{3/4},0,0)\trans,\qquad(0,0,\sqrt2,1)\trans
\]
give a real Gram completion. There is no rational matrix in the feasible fiber.

An exact exposer is
\[
 W=\diag\left(\begin{pmatrix}2&-\sqrt2\\-\sqrt2&1\end{pmatrix},
 \begin{pmatrix}\sqrt2&-2\\-2&2\sqrt2\end{pmatrix}\right).
\]
Both blocks are positive semidefinite of rank one, and $\ip{W}{A(t)}=0$ identically. In $W=W_0+\sqrt2W_1$, $\det W_0=-8$. Thus its rational zero space is $\{0\}$, and one arithmetic reduction proves infeasibility over $\Q$.

This example also rules out arbitrary positive weights in the totally real extension $\Q(\sqrt2)$: the first block equals $\sqrt2(1,\sqrt2)\trans(1,\sqrt2)$, whose weight changes sign under the second embedding. The unweighted Gram factor requires a non-totally-real extension. Finally, no separator from the original real cone can certify this arithmetic failure, because the real fiber is nonempty.
\end{example}

\section{Observation coercivity and the corrected Farkas alternatives}\label{sec:observation}
An arithmetic obstruction, ordinary strong infeasibility, and weak infeasibility must not be merged into a single strict-separation claim. For nonzero $C=C(P)$ define
\begin{equation}\label{eq:beta}
 \beta(P,\mathcal L)=\min\{\norm{\mathcal L(X)}_2:X\in C,\ \tr X=1\}.
\end{equation}

\begin{proposition}[Exact observation condition]\label{prop:beta}
The following are equivalent: $\beta>0$; $\ker\mathcal L\cap C=\{0\}$; some $\mathcal L^*y$ is strictly positive on every nonzero permitted vector; and a finite lift bound $\tr X\le\kappa\norm{\mathcal L(X)}_2$ holds on $C$. The best such $\kappa$ is $1/\beta$. Moreover
\begin{equation}\label{eq:betadual}
 \beta=\max_{\norm{y}_2\le1}\ \min_{x\in P,\ \norm{x}_2=1}x\trans\mathcal L^*(y)x.
\end{equation}
When $\beta>0$, the image $\mathcal L(C)$ is closed.
\end{proposition}
\begin{proof}
The image $B=\mathcal L(\{X\in C:\tr X=1\})$ is compact and convex. Let $z$ be its closest point to zero. If $z\ne0$, the projection inequality gives $\ip{z/\norm z}{b}\ge\norm z$ for every $b\in B$. Cauchy--Schwarz at $z$ proves the reverse inequality in the max--min formula. If $z=0$, the maximum is zero, attained by $y=0$. Minimizing a linear functional on the Gram base is the same as minimizing on normalized atoms. Homogeneity proves the equivalences and the optimal lift constant. A convergent sequence of observations has bounded trace under that constant. Since Gram matrices are positive semidefinite, their Frobenius norms are bounded by trace; closedness and a subsequence give a lift of the limit.
\end{proof}

These convex identities are classical separation consequences, included to identify the exact assumption required by a quantitative dual certificate.

\begin{proposition}[Strict rational observation certificate]\label{prop:strictsep}
Suppose $\beta>0$, $\norm b_2\le1$, and $\delta=\dist(b,\mathcal L(C))>0$. There is rational $\widehat y$ such that
\[
 \ip{\widehat y}{b}\le-3\delta/4,\qquad
 x\trans\mathcal L^*(\widehat y)x\ge\delta\beta\norm{x}_2^2/16\quad(x\in P).
\]
A rational rounding tolerance with logarithm of reciprocal
\[
 O\bigl(\log(m+1)+\log(1/\delta)+\log^+(1/\beta)+\log^+\norm{\mathcal L}\bigr)
\]
suffices. Supplied rational lower bounds for the positive quantities give explicit rational margins.
\end{proposition}
\begin{proof}
Projection onto the closed cone $\mathcal L(C)$ gives a unit $y$ in its dual with $\ip yb=-\delta$. By \eqref{eq:betadual}, choose unit $y_0$ whose adjoint has cone margin $\beta$. For $t=\delta/8$, $y+ty_0$ has target pairing at most $-7\delta/8$ and cone margin at least $\delta\beta/8$. Round it rationally in Euclidean norm with error at most
\[
 \min\{\delta/8,\ \delta\beta/[16\max(1,\norm{\mathcal L})]\}.
\]
The asserted inequalities follow from Cauchy--Schwarz and the adjoint operator norm. Coordinate rounding incurs an additional $O(\log(m+1))$ bits. The unrounded norm is bounded, so these denominator bits also bound numerator bits.
\end{proof}

\begin{example}[Two distinct failures of unconditional strict separation]\label{ex:weak}
For $P=\R_+^2$, $\mathcal L(X)=X_{11}$ and $b=-1$, the image is closed and the problem is strongly infeasible. Nevertheless every observation-supported separator is $\diag(y,0)$ and has zero margin at $e_2$. Thus a positive cone margin cannot be required without additional observability, even in a strongly infeasible problem.

For $\mathcal M(X)=(X_{11},X_{12})$, the target $(0,1)$ has no positive-semidefinite lift, but is approached by the completely positive matrices
\[
 \begin{pmatrix}\varepsilon&1\\1&1/\varepsilon\end{pmatrix},\quad\varepsilon>0.
\]
This is weak infeasibility of a nonclosed image, so no strict separating functional exists. Both examples differ from Example~\ref{ex:arithmetic}, which is real feasible.
\end{example}

\begin{corollary}[Partial Gram completion with observed diagonal]\label{cor:diagonal}
If the observation list contains every diagonal entry, then $\beta\ge1/\sqrt n$ (for the unweighted Euclidean list containing those entries). The real completion image is closed. Consequently rational partial data have the following exact trichotomy: a rational permitted Gram completion; a real completion but an arithmetic obstruction chain to rational completion; or real infeasibility with a rational observation-supported strictly positive cone separator. In each fixed dimension the alternatives are effective with the bounds of Section~\ref{sec:complexity}.
\end{corollary}
\begin{proof}
For $X\succeq0$ of trace one, $\norm{\mathcal L(X)}_2\ge\norm{\diag X}_2\ge1/\sqrt n$. Apply Theorem~\ref{thm:main} and Propositions~\ref{prop:beta}--\ref{prop:strictsep}, rescaling $b$ by a positive rational number first if its norm exceeds one.
\end{proof}

A total finite-precision alternative also allows exact side constraints. First compute the carrier for a rational affine $H_0$. If it is empty, return its arithmetic obstruction. Otherwise decide
\begin{equation}\label{eq:tolerance}
 X\in C(\AC(P,H_0))\cap H_0,\qquad\norm{\mathcal M(X)-b}_2<\rho,
\end{equation}
where $\rho>0$ is rational. If feasible, \eqref{eq:closure} and openness of the last inequality give a rational Gram completion satisfying $H_0$ exactly and the tolerance strictly. It can be constructed by moving slightly toward a terminal reserve point and repairing observations as in the next section. If infeasible, every rational completion in $H_0$ has error at least $\rho$. This disjoint, exhaustive alternative does not confuse an approximate factor with exact membership and remains valid for nonclosed observation images.


\section{Bit complexity and complete algorithmic resources}\label{sec:complexity}
The input consists of rational cone inequalities or rational ray descriptions, the rational measurement system, and, where needed, a rational target and precision. Let $L$ be the total binary input length, $s$ the total number of input cone rows after conversion to inequalities, and $g$ the number of components. Redundant measurement equations can be removed, leaving at most $d$ independent equations or an immediate inconsistency. No promise of a positive feasibility margin is made.

\begin{theorem}[Effective fixed-dimensional carrier theorem]\label{thm:complexity}
There is an absolute, effectively computable constant $c$ such that the carrier, factor, insertion, prescribed-direction and tolerance procedures above can be implemented deterministically within
\begin{equation}\label{eq:time}
 (L+p+2)^{c(n+1)^{10}}
\end{equation}
bit operations, where $p$ bounds the requested binary approximation precision and is zero for exact decisions. A conservative bound for the algebraic degrees in a reduction is $2^{c(n+1)^6}$ and for the coefficient and output coordinate bit heights is
\begin{equation}\label{eq:height}
 (L+p+1)2^{c(n+1)^6}.
\end{equation}
The carrier has a description using at most $g(s+1)^{nd}$ rational polyhedral components, each an original cone intersected with a rational subspace. There are at most $d$ exposers. A deterministic factor construction outputs at most $d\,O(h)$ rational columns when $h$ bounds the weighted-certificate height; a factor of width at most $4d$ exists with the same height order. A full-span reserve uses at most $2d$ weighted atoms before width compression.
\end{theorem}

\begin{proof}
We give the resource accounting, including the real-algebraic work. The standard sampling and one-block elimination bounds used below are those stated in Section 2 of Safey El Din--Zhi~\cite{SEDZ}, based on real-algebraic algorithms~\cite{BPR}. In $k$ variables, sampling from a formula of $J$ polynomials of degree at most $D$ has cost polynomial in coefficient bit length times $J^{k+1}D^{O(k)}$, sample degree $D^{O(k)}$, and coefficient bit height bounded by the input height times $D^{O(k)}$. Eliminating one quantified block of $a$ variables with $k$ free variables gives degree $D^{O(a)}$, a number of predicates bounded by
\[
 J^{(k+1)(a+1)}D^{(k+1)O(a)},
\]
and coefficient heights bounded by the input height times $D^{(k+1)O(a)}$. These are charged operations, not unit-cost oracles. Boolean formulas are retained as formulas rather than expanded into all combinations of component choices.

At any stage each domain is $P_i\cap U$, with $P_i$ original and $U$ rational. Equalities describing $U$ can be compressed to at most $n$ rows. Enumerate subsets $J$ of at most $n$ original inequality rows. On $U\cap\ker H_J$ choose a rational basis $B_J$, impose all coefficient kernels of $B_J\trans WB_J$, and intersect with $P_i$. Every resulting vector is a zero. The active row span of the actual minimal face of a rational zero has an independent subset of at most $n$ rows, so this enumeration captures every rational zero. Thus one update multiplies the number of components by at most $(s+1)^n$, giving the asserted component bound after at most $d$ updates. Lineality causes no change: it is handled by the rational equalities and by pairs of opposite generators.

Real feasibility is expressed by at most $d$ permitted columns:
\[
 \exists u_1,\ldots,u_d:\quad u_j\in P,\quad
 \mathcal L\Bigl(\sum_ju_ju_j\trans\Bigr)=b.
\]
There are at most $nd$ variables and degree two. Relative strict feasibility adds $\tau>0$ and $\tau B$ to that sum, using at most $nd+1$ variables and degree two. Sampling or infeasibility of these formulas has the stated bounded-dimensional complexity.

An exposer has $d$ free matrix variables and one universal block of $n$ vector variables. Nonnegativity on the domain is a Boolean implication whose polynomial degree is at most three; vanishing on $H\cap V$ and $\ip WB=1$ are rational linear equations. Apply one-block elimination and then sampling in at most $d$ variables. The sampled matrix lies in a number field of degree $2^{O((n+1)^3)}$ with coefficient height bounded by the current rational height times $2^{O((n+1)^3)}$. Factor its univariate representation over $\Q$, select the isolating root, invert its nonzero denominators modulo the irreducible polynomial, and expand in a power basis. Rational kernel computation then discards the algebraic field entirely. Hence successive steps multiply heights by dimension-dependent factors; they do not form a tower of field extensions. At most $d$ such factors, together with rational Cramer bounds for subspace bases and reserves, are absorbed by \eqref{eq:height}.

For clarity about the time exponent, the component count makes the number of original predicates at most $(L+2)^{O((n+1)^3)}$. One-block elimination with $k=d$, $a=n$ raises this by $O((n+1)^3)$; sampling in $d$ variables raises it by a further $O((n+1)^2)$. This is bounded by $(L+2)^{O((n+1)^8)}$. Rational row operations, field factorization, at most $d$ repetitions, and the terminal constructions below are absorbed by the conservative exponent in \eqref{eq:time}. Conversion from a rational ray description in fixed dimension enumerates at most a power $O(n)$ of its input rays and is absorbed as well. Enlarging one absolute constant $c$ covers all primitive constants. The displayed exponents are not claimed sharp.

To construct a terminal rational completion, sample algebraic $\tau>0$ and permitted $u_j$ satisfying
\[
 \mathcal L\Bigl(\sum_ju_ju_j\trans+\tau B\Bigr)=b.
\]
Write each $u_j$ as a nonnegative combination of rational generators of a containing component, and approximate those coefficients by nonnegative rationals to obtain $r_j$ within that component. At most $n$ generators per column suffice; lineality is represented using both signs. Let $M$ have columns $\mathcal L(G_i)$ and choose a rational linear right inverse $R$ on $\operatorname{im}M$. For a positive rational approximation $\tau_0$ to $\tau$, define
\begin{equation}\label{eq:repair}
 e=\tau_0\mathbf1+R\left[b-\mathcal L\Bigl(\sum_jr_jr_j\trans\Bigr)-\tau_0M\mathbf1\right].
\end{equation}
The bracket belongs to $\operatorname{im}M$. Thus $X_0=\sum_jr_jr_j\trans+\sum_ie_iG_i$ has the exact observations. As the approximations improve, $e\to\tau\mathbf1$, so all entries become positive. Algebraic root separation bounds $\tau$ below by $2^{-h2^{O((n+1)^3)}}$ for current height $h$; Cramer's rule bounds $R$ and the generator coordinates. Consequently a number of dyadic refinement bits bounded by \eqref{eq:height} suffices. Exact positivity checks determine success. When $M=0$, feasibility gives $b=0$, and one simply takes $R=0$.

For a rational target in the terminal closed slice, take a sufficiently small rational segment toward $X_0$ to meet $2^{-p}$ accuracy, then apply the same construction to the rational segment point. For \eqref{eq:tolerance}, an algebraic feasible sample with a strict error slack can first be moved slightly toward an algebraic reserve completion; root separation bounds the required positive slack and interpolation parameter. Formula \eqref{eq:repair} then preserves the exact observations without destroying the open tolerance. Prescribed-direction tests are one-dimensional real feasibility tests after the same reduction. All coefficients obey \eqref{eq:height} with $L+p$ in place of $L$.

Finally, after rational dependence elimination there are at most $d$ positive rational weights $a/b$ of height $h$. Four squares prove existence of at most $4d$ Gram columns, with no larger height order. For a deterministic polynomial construction without a four-square subroutine, expand the positive integer $ab$ in binary. A term $2^{2j}$ is one square and a term $2^{2j+1}$ is two equal squares. Dividing their square roots by $b$ gives $O(h)$ rational columns per weighted atom. Their signs are chosen nonnegative, so each column remains in its permitted cone. This proves the resource distinctions in the statement.
\end{proof}

The exact arithmetic obstruction certificate consists of the algebraic exposers with isolating data, their nonnegativity/elimination certificates, rational coefficient-kernel restrictions, and the final real infeasibility certificate. A reproducible transcript of the specified exact decision and elimination computations has length bounded by their charged running time; it need not be a full cylindrical cell decomposition. There is no assertion of a short variable-dimensional coNP certificate. Similarly, the procedure is polynomial for each fixed $n$, not uniformly polynomial in $n$.

If the observation dimension is large, Euclidean norm constraints are still evaluated in at most $d$ matrix coordinates. Dual observation vectors can be restricted to the rational span of the columns of $\mathcal L$ and $b$, of dimension at most $d+1$; its Gram matrix retains the Euclidean norm exactly. Thus the observation count is part of $L$, not an uncharged extra fixed-dimension assumption.


\section{Sharp angular corrections for bounded Gram moments}\label{sec:angular}
The carrier theorem is constructive even on a boundary fiber. A particularly explicit application is the genuinely nonsimplicial cone
\[
 K_r=\{(t,z)\in\R^{r+1}:t\ge0,\ |z_i|\le t\},\qquad r\ge2.
\]
Its rank-one Gram matrices describe unnormalized moments of bounded variables. We preserve mass, all first moments and all cross moments, and perturb only the individual second moments.

Let $G=\{-1=a_0<a_1<\cdots<a_m=1\}$ be a scalar grid. Define
\begin{align}
 e(G)&=\max_{0\le j<m}\frac{1-a_ja_{j+1}-\sqrt{(1-a_j^2)(1-a_{j+1}^2)}}2\label{eq:eg}\\
 &=\max_{0\le j<m}\sin^2\left(\frac{\arccos a_j-\arccos a_{j+1}}2\right).
\end{align}
For possibly different grids $G_i$, let $C_G$ be generated by $(1,z)(1,z)\trans$ with $z_i\in G_i$. Define $\Phi_e$ to preserve $A_{00}$ and every off-diagonal entry and set
\begin{equation}\label{eq:phie}
 (\Phi_e A)_{ii}=(1-e_i)A_{ii}+e_iA_{00}\quad(1\le i\le r).
\end{equation}

\begin{theorem}[Exact grid correction]\label{thm:angular}
The inclusion $\Phi_e(C(K_r))\subseteq C_G$ holds if and only if
\[
 e(G_i)\le e_i\le1\qquad(1\le i\le r).
\]
The corresponding dual inclusion is $\Phi_e^*(C_G^*)\subseteq\COP(K_r)$. The adjoint adds $\sum_i e_iQ_{ii}$ to $Q_{00}$, multiplies $Q_{ii}$ by $1-e_i$, and preserves off-diagonals. For rational grids every coefficient below the bound has a strict rational rank-one counterexample and a rational grid-positive quadratic witness.
\end{theorem}
\begin{proof}
First consider one coordinate with mean $y\in[a,b]$, where $a,b$ are consecutive grid nodes. For a grid-valued random variable $Z$, $(Z-a)(Z-b)\ge0$, so
\[
 \operatorname{Var}Z\ge(y-a)(b-y).
\]
The distribution on $a,b$ of mean $y$ attains equality. Since $Z^2\le1$, the largest variance is $1-y^2$, attained by the distribution on $-1,1$ of mean $y$. Mixing these two laws realizes every variance between the two bounds. Therefore the desired second moment $y^2+e(1-y^2)$ is possible precisely when
\[
 e\ge\sup_{a<y<b}\frac{(y-a)(b-y)}{1-y^2},\qquad e\le1.
\]
For interior endpoints, nonnegativity of
\[
 e(1-y^2)-(y-a)(b-y)=(1-e)y^2-(a+b)y+ab+e
\]
has its threshold at the smaller discriminant root
\[
 e=\frac{1-ab-\sqrt{(1-a^2)(1-b^2)}}2.
\]
The double root lies in $[a,b]$. For instance, write $a=\cos\alpha$, $b=\cos\beta$, $\alpha>\beta$; it equals $\cos((\alpha+\beta)/2)/\cos((\alpha-\beta)/2)$, which lies between $a$ and $b$. When $a=-1$ or $b=1$, the same value follows by a one-sided limit; the whole interval $[-1,1]$ has threshold one.

Choose the coordinate laws independently. They preserve each mean and every cross product, while producing exactly \eqref{eq:phie} on a normalized rank-one atom. Homogenize by $t^2$ and sum over real permitted atoms to obtain the inclusion. If an $e_i>1$, the corrected atom with mean zero has second moment greater than its mass and is not even in $C(K_r)$.

For the remaining necessity, the rational form
\[
 Q(t,z)=(z_i-at)(z_i-bt)
\]
is nonnegative at every grid atom. Its pairing with the corrected atom of mean $y$ is
\[
 -(y-a)(b-y)+e_i(1-y^2).
\]
Below the threshold this is strictly negative for some $y$ in the gap; for rational nodes a rational such $y$ exists by continuity. Duality yields the adjoint statement, whose formula follows by taking the trace pairing. It is important that $\Phi_e$ is generally not self-adjoint.
\end{proof}

\begin{theorem}[Optimal grids and their arithmetic obstruction]\label{thm:optimalgrid}
Among all grids with $m$ intervals the least correction is
\[
 \sin^2\bigl(\pi/(2m)\bigr),
\]
attained only by equally spaced angular nodes. Among rational grids this is an infimum, attained exactly for $m=1,2,3$. The explicit rational grids
\begin{equation}\label{eq:stereographic}
 G_q=\left\{\pm\frac{q^2-j^2}{q^2+j^2}:0\le j\le q\right\},\qquad q\ge1,
\end{equation}
have $2q+1$ nodes, coordinate height $O(\log(q+1))$, and exact correction
\begin{equation}\label{eq:rationalrate}
 e(G_q)=\frac1{q^2+1}.
\end{equation}
Their correction is at most $16/\pi^2$ times the geometric optimum for $2q$ intervals. By contrast, the equally spaced scalar grid with $m$ intervals has exact correction $1/m$.
\end{theorem}
\begin{proof}
The angular gaps sum to $\pi$, hence their maximum is at least $\pi/m$. Equality forces all gaps equal. The resulting nodes are the cosines of multiples of $\pi/m$. These are rational for $m=1,2,3$. For $m\ge4$, rationality would make $2\cos(\pi/m)$ a rational algebraic integer strictly between one and two, which is impossible. Rational grids approximate the optimal nodes arbitrarily closely while preserving their order, proving the infimum assertion.

The nonnegative half of \eqref{eq:stereographic} has angles $2\arctan(j/q)$ from zero to $\pi/2$. Since $\arctan$ is increasing and concave on $[0,1]$, the largest angular increment is $2\arctan(1/q)$. Reflection gives the same bound on the other half. Formula \eqref{eq:eg} gives \eqref{eq:rationalrate}. The rational numerator and denominator are at most $2q^2$.

For $x=\pi/(4q)$, $\tan x\ge x$ implies $\sin^2x\ge x^2/(1+x^2)$. Consequently
\[
 \frac{1/(q^2+1)}{\sin^2(\pi/(4q))}\le
 \frac{16q^2+\pi^2}{\pi^2(q^2+1)}\le\frac{16}{\pi^2}.
\]
For the equally spaced grid, the first and last angular gaps are the largest: the derivative of $\arccos$ is symmetric and increases in magnitude toward the endpoints, so its integral over an interval of fixed length is largest there. The endpoint formula in \eqref{eq:eg} gives $1/m$.
\end{proof}

\begin{corollary}[Exact bounded-moment rationalization]\label{cor:moments}
For rational $A\in C(K_r)$ put
\[
 D=\diag(0,A_{00}-A_{11},\ldots,A_{00}-A_{rr}).
\]
For every rational $0<\varepsilon\le1$, $A+\varepsilon D$ has rational permitted Gram columns, preserving mass, first moments, all cross moments, and every saturated second moment. A grid with $q^2+1\ge1/\varepsilon$ suffices. There is a representation using at most $d=(r+1)(r+2)/2$ weighted rational atoms with heights polynomial in the binary encoding length of $(A,\varepsilon)$ for fixed $r$, and hence at most $4d$ rational columns with the same height order.
\end{corollary}
\begin{proof}
Theorem~\ref{thm:angular} puts the corrected rational matrix in a finite rational cone. A rational basic feasible representation uses at most $d$ independent atoms. For a grid atom, multiply by the product of its coordinate denominators, at most $(2q^2)^r$, to get an integer vector. Cramer's rule on at most $d$ independent integer atom columns bounds the rational weights by
\[
 O\bigl(L_\varepsilon+d(r\log(q+1)+\log(d+1))\bigr)
\]
bits, where $L_\varepsilon$ is the total encoding of the corrected target. Lemma~\ref{lem:reserve} proves the width assertion. If $A_{ii}=A_{00}$ the correction is zero in that coordinate; probabilistically every original atom is already on one of the two corresponding facets. Thus saturated faces are retained exactly.
\end{proof}

For a rational box $l_i\le z_i/t\le u_i$, affine homogenization gives the exact correction
\[
 A_{ii}\longmapsto(1-e_i)A_{ii}+e_i\bigl((l_i+u_i)A_{0i}-l_iu_iA_{00}\bigr).
\]
The slack being changed is the expectation of $(z_i-l_it)(u_it-z_i)$. This is a constructive carrier certificate for an entire non-simplicial family of measured Gram fibers, not a congruence of a simplicial orthant. Enumerating $(2q+1)^r$ atoms is nevertheless exponential in binary accuracy. The fixed-dimensional nonlocal algorithm of Theorem~\ref{thm:complexity} is a different resource and is not replaced by that enumeration.

Grid hierarchies and quadratic approximation rates predate this work~\cite{Yildirim,DKLSV}. The assertions isolated here are the exact measurement-compatible angular constant, its grid optimization and rational nonattainment, and the explicit rational stereographic construction.


\section{Sharp lattice conditioning and what it controls}\label{sec:lattice}
We record the quantitative lattice estimate needed to distinguish coefficient growth from traversal. The orthant estimate and its sharp shear construction were established in the preceding Manuscript II checkpoints~\cite{M2prior}; the qualitative Interior Ryshkov property and the underlying Diophantine method are prior literature~\cite{DSV21,OS26}.

\begin{lemma}[Orthant norm-to-margin bound]\label{lem:orthantmargin}
If $P[v]\ge1$ for every $0\ne v\in\Z_{\ge0}^s$ and $M\ge\max(1,\norm P_F)$, then
\[
 \min_{x\ge0,\ \norm x_2=1}P[x]\ge\frac1{8s(16sM)^{s-1}}.
\]
\end{lemma}
\begin{proof}
A negative value at a real nonnegative vector would persist at a rational vector and contradict the integer inequalities after clearing denominators. Thus $P$ is copositive. Choose a unit minimizer $x$ with value $\mu$, and let its positive support have size $a$. Stationarity on the support gives $P_{SS}x_S=\mu x_S$. If $a=1$, the claim follows from $P_{ii}\ge1$. Otherwise normalize a largest coordinate to one, writing $y=x/x_j$ with $\norm y_2\le\sqrt s$.

Let $Q=\lceil\sqrt{8sM}\rceil$. The pigeonhole proof of simultaneous Dirichlet approximation gives $1\le m\le Q^{a-1}$ and an integer nonnegative vector $v$ with $v_j=m$ and $|v_i-my_i|\le1/Q$ on the other support coordinates. Nonnegativity follows since the approximated numbers are nonnegative and the error is less than one. Put $e=v-my$ and use support stationarity to get
\[
 1\le P[v]\le\mu(m\norm y+\norm e)^2+M\norm e^2
 \le4s\mu Q^{2(s-1)}+1/8.
\]
Since $Q\le4\sqrt{sM}$, the stated weaker constant follows. The argument proves strict positivity rather than presupposing it.
\end{proof}

\begin{theorem}[Cone-relative conditioning, with sharp exponent]\label{thm:lattice}
Let a rational polyhedral cone $K$ of span dimension $r$ be generated by integer vectors of norm at most $R\ge1$. If $P[v]\ge1$ for every nonzero integer $v\in K$ and $M\ge\max(1,\norm P_F)$, then
\begin{equation}\label{eq:Kmargin}
 \min_{x\in K,\ \norm x_2=1}P[x]\ge
 \frac1{8r^2R^2(16r^2R^2M)^{r-1}}.
\end{equation}
For every fixed rational polyhedral cone with relative dimension $r\ge2$, the exponent $r-1$ cannot be reduced, even on relatively perfect positive-definite forms.
\end{theorem}
\begin{proof}
Every $x\in K$ is a nonnegative combination $Ba$ of $s\le r$ independent integer generators. The pullback $B\trans PB$ has nonnegative integer minimum at least one, norm at most $sR^2M$, and $\norm{Ba}^2\le sR^2\norm a^2$. Apply Lemma~\ref{lem:orthantmargin} and enlarge constants from $s$ to $r$ to get \eqref{eq:Kmargin}.

For sharpness, use the relative lattice $\lin K\cap\Z^n$. It has a basis $B$ lying in $K$: choose a primitive relative-interior lattice vector, extend it to a lattice basis, and add sufficiently large integer multiples of it to the other basis vectors. The relative-interior vector can be made primitive by dividing its lattice gcd. Let $T_r$ have diagonal one and adjacent off-diagonal $-1/2$, and let $V_q=I-qS$ for the upper shift $S$ and integer $q\ge2$. Define the form by
\[
 P_q[Bz]=z\trans V_q\trans T_rV_qz.
\]
For every nonzero integer $z$,
\[
 T_r[z]=\tfrac12\left(z_1^2+z_r^2+\sum_{i=1}^{r-1}(z_i-z_{i+1})^2\right)\ge1.
\]
The positive interval indicators attain one and their squares span $\Ssym^r$. Since $V_q^{-1}=I+qS+\cdots+q^{r-1}S^{r-1}$ is nonnegative integral, their inverse images under $V_q$, then under $B$, lie in $K$ and prove relative perfection. The norm of $P_q$ is $O_K(q^2)$. The minimal vector $BV_q^{-1}e_r$ has norm $\Omega_K(q^{r-1})$ and value one, so the margin is $O_K(q^{-2(r-1)})$. No smaller norm exponent is possible.
\end{proof}

The bound limits minimal vector lengths. In full dimension, writing its reciprocal as $U$, every active integer vector has squared norm at most $U$. In upper-triangular coordinates its integer constraint normal has entries at most $2U$. Cramer's rule bounds vertex entries and primitive tangent-ray entries by $d![2\lceil U\rceil]^d$. Thus coefficient heights grow only logarithmically with $M$ for fixed dimension. These bounds do not count how many normal-fan cells are visited. The next theorem proves that this missing resource is genuinely independent of the arithmetic carrier.

\section{An exponential local walk inside one arithmetic face}\label{sec:paths}
Fix $r\ge2$, put $n=r+2$, and consider
\begin{equation}\label{eq:pathcone}
 K=\{(t,z_1,\ldots,z_r,y):t,y\ge0,\ |z_i|\le t\}.
\end{equation}
This is the product of a homogenized cube and a ray. It has $2^r+1>n$ extreme rays and is not simplicial. Its generalized Ryshkov set consists of $P$ with $P[v]\ge1$ for all nonzero $v\in K\cap\Z^n$. At a perfect vertex let the shadow rule use the fixed anchor $C_0$ and minimize $\ip{C_0}{R}/(-\ip A R)$ among objective-descending extreme tangent rays. We count exact adjacent moves, not arbitrary preprocessing or nonlocal factor recognition.

\begin{theorem}[Depth-zero local/global separation]\label{thm:path}
For the fixed cone \eqref{eq:pathcone} and every integer $N\ge2$, there is a rational interior target $A_N$ with the following properties.
\begin{enumerate}[label=(\roman*)]
\item Its input length is $O_n(\log N)$. It has at most $6\cdot3^r$ explicit rational Gram columns, each of height $O_n(\log N)$.
\item Its arithmetic carrier is the whole $K$, and its arithmetic reduction depth is zero. For full observation, $\beta(K,\mathrm{id})=1/\sqrt n$, independent of $N$.
\item From a fixed perfect vertex and its full-minimum anchor, the specified shadow rule has a unique choice at every crossing and makes exactly $N-1$ adjacent pivots.
\item After rational-square normalization, $\norm{A_N}_F\le1$ and its interior distance is at least
\begin{equation}\label{eq:pathradius}
 \frac{2^{-(2r+4)}}{n^2d!(N+1)^{2d}}.
\end{equation}
Every visited vertex, ray and compressed certificate has height $O_n(\log N)$.
\end{enumerate}
A fixed integral unimodular congruence makes all target entries strictly positive, preserving the local path and the conclusions up to fixed-dimensional constants.
\end{theorem}

\begin{proof}
For $k\ge1$ define a quadratic form, independent of $z$, by
\begin{equation}\label{eq:Pk}
 P_k=t^2-(2k+1)ty+(k^2+k+1)y^2.
\end{equation}
With $a=t-ky$ this is $a^2-ay+y^2$. Its allowed nonzero integer minima are exactly the vectors with
\begin{equation}\label{eq:minima}
 (t,y)=(1,0),(k,1),(k+1,1),\qquad z\in\Z^r,\quad |z_i|\le t.
\end{equation}
Indeed the binary integer form has minimum one; if $y\ge2$, completing the square gives value at least $3y^2/4\ge3$. The cases $y=0,1$ give precisely the stated possibilities. If $t=y=0$, cone membership forces $z=0$.

These minima prove perfection. A quadratic form vanishing on $(1,z,0)$ for all $z\in\{-1,0,1\}^r$ has zero $tt,tz,zz$ coefficients by interpolation on three nodes in each coordinate. Vanishing on $(k,z,1)$ next makes every $zy$ coefficient zero and imposes $2kQ_{ty}+Q_{yy}=0$. The family with $k+1$ gives $Q_{ty}=Q_{yy}=0$. Thus the active normals span $\Ssym^n$.

For $s\ge2$ the common active families at $P_{s-1},P_s$ have base coordinates $(1,0),(s,1)$. The same calculation leaves the one-dimensional annihilator
\[
 R_{s-1}=P_s-P_{s-1},\qquad (R_{s-1})_{ty}=-1,\quad(R_{s-1})_{yy}=2s.
\]
The segment between the two feasible vertices is an edge. Its length in this normalization is exactly one: a vector with base $(s+1,1)$ becomes infeasible immediately beyond $P_s$. Relative strict copositivity of $P_k$ on $K$ follows from the positive-definite binary form and $\norm z^2\le r t^2$, so the local normal cones are generated by the finitely many active vectors.

At $P_1$ the anchor obtained by summing \emph{all} minima has binary entries
\begin{equation}\label{eq:anchor}
 a=2\cdot3^r+4\cdot5^r,\quad b=3^r+2\cdot5^r,\quad c=3^r+5^r
\end{equation}
and each $z$ diagonal equal to $\zeta=4\cdot3^{r-1}+10\cdot5^{r-1}$. All other $z$ entries vanish. Note $a=2b$, $2c-b=3^r$, and $0<\zeta<a$.

Write $T=N+1/2$ and $U=N^2+N+1$. Define the unnormalized target by its binary block
\begin{equation}\label{eq:target}
 A_N|_{t,y}=\begin{pmatrix}U&T\\T&1\end{pmatrix},\qquad
 A_N|_z=(U/2)I_r,
\end{equation}
with all remaining entries zero. Its binary determinant is $3/4$, hence it is positive definite. An explicit permitted factorization starts with base vectors $(1,0),(N,1),(N+1,1)$, each of weight $1/2$, and independently assigns $z_i=-t,0,t$ with probabilities $1/4,1/2,1/4$. The $3\cdot3^r$ resulting weighted atoms have Gram sum \eqref{eq:target}. All weights are positive dyadic numbers and their normals span $\Ssym^n$ by the preceding interpolation argument. Thus $A_N\in\operatorname{int}C(K)$. Each dyadic weight is one rational square or twice one rational square, giving at most $6\cdot3^r$ Gram columns.

For $2\le s\le N$ set
\begin{equation}\label{eq:lambda}
 \lambda_s=\frac{sc-b}{T-s},\qquad C_s=C_0+\lambda_s A_N.
\end{equation}
The parameters are positive and strictly increasing, because $2c-b>0$ and $cT-b>0$. The binary part of $C_s$ has the exact decomposition
\begin{equation}\label{eq:binarynormal}
 \chi_s(1,0)(1,0)\trans+(c+\lambda_s)(s,1)(s,1)\trans,
\end{equation}
where
\begin{align*}
 \chi_s&=a-sb+\lambda_s(U-sT)\\
 &=a-bT+s(cT-b)+\tfrac34\lambda_s>0.
\end{align*}
For positivity, $s\ge2$ and $cT-b>0$ make the last expression at least $(2c-b)T+\tfrac34\lambda_s$.

Let
\[
 v_s=\frac{\zeta+\lambda_s U/2}{a+\lambda_s U}\in(0,1).
\]
Lift each base vector in \eqref{eq:binarynormal} independently by $z_i=-t,0,t$ with probabilities $v_s/2,1-v_s,v_s/2$. This gives exactly $C_s$, with strictly positive weights on normals spanning the common active hyperplane. Hence $C_s$ lies in the relative interior of the shared normal-cone facet. Its sole zero tangent ray is $R_{s-1}$, and
\[
 \ip{A_N}{R_{s-1}}=-2(T-s)<0.
\]
Between two successive crossings the positive combination of the two facet certificates uses the three base families $(1,0),(s,1),(s+1,1)$, whose normals span the full matrix space. Therefore the intervening objective is in the interior normal cone at $P_s$. Initially $C_0$ is interior at $P_1$. Finally $A_N$ itself is interior in the normal cone at $P_N$, so every objective after the last crossing remains there. This proves the exact unique shadow trajectory $P_1,P_2,\ldots,P_N$ and its $N-1$ pivots. No chosen ray can be a copositive separator: each has a negative permitted value, and the target is completely positive. The statement does not prohibit a different algorithm from recognizing the displayed factorization immediately.

For the interior radius, select $d$ independent integer atoms from the displayed target representation. Their upper-triangle coordinate matrix has integer determinant of absolute value at least one and entries bounded by $(N+1)^2$. Its inverse entries have magnitude at most $(d-1)!(N+1)^{2(d-1)}$. Every target weight is at least $w_0=2^{-(2r+1)}$. A Frobenius perturbation of norm at most
\[
 \eta=\frac{w_0}{2d!(N+1)^{2(d-1)}}
\]
can consequently be absorbed into those $d$ coefficients while keeping them positive. Normalize by the rational square $[2n(N+1)]^2$. This bounds the Frobenius norm by one and gives \eqref{eq:pathradius}. All expressions involve rational numbers with $O_n(\log N)$ bits, including $s\le N$ in the compressed normal certificates.

The carrier is $K$ because the target is an interior rational point; Theorem~\ref{thm:main} or the reserve argument applies with no reduction. For full observation the trace-one PSD norm is at least $1/\sqrt n$. Moreover $I\in C(K)$: independent signs in $(1,z,0)$ generate $\diag(1,I_r,0)$, and the $y$ axis supplies the last diagonal. Hence $\beta(K,\mathrm{id})=1/\sqrt n$.

Finally apply the fixed integral map
\[
 (t,z,y)\longmapsto(t,z+2t\mathbf1,y+t).
\]
Its determinant is one. It sends $K$ to the nonsimplicial cone $t\ge0$, $t\le z_i'\le3t$, $y'\ge t$. The congruenced target is positive definite and every entry is strictly positive. Integral minima, full-minimum anchors, trace pairings, extreme rays and the minimum-ratio rule are carried bijectively by the inverse-transpose action on forms. Norm and interior-radius estimates change only by constants depending on $n$. This finishes the proof.
\end{proof}

\begin{corollary}[No polynomial walk bound from genuine arithmetic faces]\label{cor:barrier}
For the specified fixed-start rule, no polynomial in the binary input length, rational factor height and width, binary interior condition, and arithmetic reduction depth bounds the number of adjacent pivots, even after fixing the vector cone, its arithmetic carrier, and the observation coercivity. The entire shadow objective line can remain in the same genuine interior matrix-cone face while crossing exponentially many perfect-form cells.
\end{corollary}
\begin{proof}
Take $N=2^B$ in Theorem~\ref{thm:path}. Every listed binary resource is $O_n(B)$ or fixed, the depth is zero, and the pivot count is $2^B-1$. The anchor and target are interior, hence so is their positive objective combination throughout the path.
\end{proof}

This is an obstruction to a particular proposed complexity inference, not a claim that no imaginable invariant can encode path length, and not a lower bound for every descending local rule. The earlier Manuscript II results provide complementary nonmember instances~\cite{M2prior}; they are not needed for this interior argument. The shear/perfectness architecture has antecedents in Dannenberg--Sch\"urmann~\cite{DS23}. What is isolated here is the fixed nonsimplicial, depth-zero, uniquely forced path with an explicit short rational certificate and inverse-polynomial interior radius.


\section{Verification, provenance, and limitations}\label{sec:verification}
The release supplies separate standard-library generator and verifier programs. The verifier imports no generator code. It checks five angular-grid cases, eleven exact facet certificates, the four-dimensional arithmetic obstruction, full rank of the target normal sets, rank $d-1$ of the facet normal sets, rational reconstruction identities and strict signs. It rejects three deliberately corrupted data sets. One target uses $N=2^{64}$; its universal proof gives $2^{64}-1$ pivots. The program checks compressed certificates, not that many simulated steps.

These tests are independent implementations of finite identities, not independent human review, a proof-assistant formalization, or a production implementation of general real quantifier elimination. The all-dimensional conclusions rest on the proofs above. The release records exact data hashes and deterministic regeneration commands.

Several distinctions are essential to the contribution. Qualitative generalized perfectness is already available in 2026~\cite{OS26}. Rational interior factorization and classical ellipsoid/separation algorithms are not new. Ordinary real facial reduction does not by itself determine rational admissible columns; the coefficient-kernel restriction and carrier canonicity are the arithmetic statements. The totally real result is a constrained extension, not a claim of first trace descent. Grid convergence rates are old; the angular optimum is specific to the stated moment-preserving map. The sharp lattice exponent is inherited from the preceding project and proved here for completeness. Finally, the local lower bound concerns the specified shadow rule and cannot be advertised as a theorem about all algorithms.

Algebraic zero and kernel restrictions are also explicit antecedents in Laplagne's exact sum-of-squares facial reduction~\cite{Laplagne}. His Section 3.4 cautions that imposing rational coefficient conditions on a general algebraic parametrization need not preserve every rational solution. Here the coordinates of a rational column in a rational face basis are themselves rational; coefficient independence therefore gives necessary restrictions. The claim isolated by Theorem~\ref{thm:main} is exact columnwise preservation and a canonical finite polyhedral carrier, not the first use of algebraic coefficients or trace kernels.

The current literature audit found no source establishing the full canonical arithmetic-carrier theorem with these exact column constraints, but a negative search is not a historical priority proof. Reported 2026 boundary CP counterexamples are not assumed in any argument, and their complete external certificates have not been independently rerun for this draft. The registered article-level venue comparison is supplied separately; no editorial probability or upper-quintile certification is implied by the mathematics or by this manuscript's length.

\begin{thebibliography}{99}
\bibitem{BW} J. M. Borwein and H. Wolkowicz, Facial reduction for a cone-convex programming problem, \emph{J. Austral. Math. Soc. Ser. A} 30 (1981), 369--380. DOI: 10.1017/S1446788700017250.
\bibitem{BPR} S. Basu, R. Pollack and M.-F. Roy, \emph{Algorithms in Real Algebraic Geometry}, second edition, Springer, 2006. DOI: 10.1007/3-540-33099-2.
\bibitem{SEDZ} M. Safey El Din and L. Zhi, Computing rational points in convex semi-algebraic sets and sum of squares decompositions, \emph{SIAM J. Optim.} 20 (2010), 2876--2889. DOI: 10.1137/090772459.
\bibitem{DSV17} M. Dutour Sikiri\'c, A. Sch\"urmann and F. Vallentin, Rational factorizations of completely positive matrices, \emph{Linear Algebra Appl.} 523 (2017), 46--51. DOI: 10.1016/j.laa.2017.02.017.
\bibitem{DSV21} M. Dutour Sikiri\'c, A. Sch\"urmann and F. Vallentin, A simplex algorithm for rational cp-factorization, \emph{Math. Program.} 187 (2021), 25--45. DOI: 10.1007/s10107-020-01467-4.
\bibitem{DS23} V. Dannenberg and A. Sch\"urmann, Perfect copositive matrices, \emph{Commun. Math.} 31 (2023), 137--156. DOI: 10.46298/cm.11141.
\bibitem{OS26} A. Oertel and A. Sch\"urmann, Generalized Perfect Matrices, arXiv:2602.05841v2 (2026).
\bibitem{Hillar} C. J. Hillar, Sums of squares over totally real fields are rational sums of squares, \emph{Proc. Amer. Math. Soc.} 137 (2009), 921--930. DOI: 10.1090/S0002-9939-08-09641-X.
\bibitem{Scheiderer} C. Scheiderer, Sums of squares of polynomials with rational coefficients, \emph{J. Eur. Math. Soc.} 18 (2016), 1495--1513. DOI: 10.4171/JEMS/620.
\bibitem{Laplagne} S. Laplagne, Facial reduction for exact polynomial sum of squares decomposition, \emph{Math. Comp.} 89 (2020), 859--877. DOI: 10.1090/mcom/3476.
\bibitem{LMT} B. F. Louren\c co, M. Muramatsu and T. Tsuchiya, Solving SDP completely with an interior point oracle, \emph{Optim. Methods Softw.} 36 (2021), 425--471. DOI: 10.1080/10556788.2020.1850720.
\bibitem{SDB} N. Shaked-Monderer, M. D\"ur and A. Berman, Complete positivity over the rationals, \emph{Pure Appl. Funct. Anal.} 3 (2018), 681--691.
\bibitem{Kostyukova} O. I. Kostyukova, Non-exposed polyhedral faces of the completely positive cone, \emph{Linear Multilinear Algebra} 73 (2025), 368--395. DOI: 10.1080/03081087.2024.2346313.
\bibitem{NL26} M. Nishijima and B. F. Louren\c co, Facial structure of copositive and completely positive cones over a second-order cone, \emph{Linear Algebra Appl.} 743 (2026), 211--249. DOI: 10.1016/j.laa.2026.04.019.
\bibitem{Yildirim} E. A. Y\i ld\i r\i m, On the accuracy of uniform polyhedral approximations of the copositive cone, \emph{Optim. Methods Softw.} 27 (2012), 155--173. DOI: 10.1080/10556788.2010.540014.
\bibitem{DKLSV} E. de Klerk, M. Laurent, Z. Sun and J. C. Vera, On the convergence rate of grid search for polynomial optimization over the simplex, \emph{Optim. Lett.} 11 (2017), 597--608. DOI: 10.1007/s11590-016-1023-7.
\bibitem{M2prior} LAA Paper project, Manuscript II quantitative checkpoints, immutable repository snapshot 98740002955f5bd9763137fc583a8d897b136ffe, \texttt{research/pro\_m2\_quantitative\_20261114}, 14 November 2026. Internal antecedent, not a published priority claim.
\end{thebibliography}


\end{document}
